\documentclass[11pt]{article}
\usepackage[utf8]{inputenc}
\usepackage{amsmath, amssymb, amsthm}
\usepackage{geometry}
\geometry{a4paper, margin=1in}
\usepackage{hyperref}
\usepackage{enumitem}

\title{Variable-Coefficient Equations: Cauchy-Euler Problem Sets}
\author{Comprehensive Practice Document}
\date{}

\begin{document}

\maketitle

\tableofcontents
\newpage

% ---------------------------------------------------------
\section{Concept Checklist \& Question Mapping}
\label{sec:mapping_start}
This section outlines all the core concepts and problem types extracted from the provided text on Variable-Coefficient Equations (specifically Cauchy-Euler equations). Each concept is mapped to its corresponding problems in the Problem Sets section to ensure comprehensive coverage.

\begin{itemize}
    \item \textbf{Concept 1: Existence and Uniqueness Analysis} 
    \textit{Identifying intervals of validity by converting to standard form and finding domains of continuity.}
    \textbf{Questions: 1, 2, 3}
    
    \item \textbf{Concept 2: Homogeneous Cauchy-Euler (Distinct Real Roots)} 
    \textit{Solving $at^2y'' + bty' + cy = 0$ where the characteristic equation $ar^2 + (b-a)r + c = 0$ yields $r_1 \neq r_2$.}
    \textbf{Questions: 4, 5, 6}
    
    \item \textbf{Concept 3: Homogeneous Cauchy-Euler (Repeated Real Roots)} 
    \textit{Handling roots where $r_1 = r_2 = r$, requiring the second linearly independent solution to be multiplied by $\ln t$.}
    \textbf{Questions: 7, 8, 9}
    
    \item \textbf{Concept 4: Homogeneous Cauchy-Euler (Complex Roots)} 
    \textit{Handling roots $r = \alpha \pm i\beta$, utilizing Euler's formula to yield $t^\alpha(c_1 \cos(\beta \ln t) + c_2 \sin(\beta \ln t))$.}
    \textbf{Questions: 10, 11, 12}
    
    \item \textbf{Concept 5: Domain Adaptation ($t < 0$)} 
    \textit{Adapting solutions for negative domains by replacing $t$ with $-t$ inside the natural logarithm to maintain real-valued functions.}
    \textbf{Questions: 13, 14, 15}
    
    \item \textbf{Concept 6: Initial Value Problems (IVPs)} 
    \textit{Applying initial conditions $y(t_0)$ and $y'(t_0)$ to find specific constants $c_1$ and $c_2$.}
    \textbf{Questions: 16, 17, 18}
    
    \item \textbf{Concept 7: Nonhomogeneous Equations (Variation of Parameters)} 
    \textit{Solving $at^2y'' + bty' + cy = f(t)$ by first converting to standard form to identify $g(t)$, then computing the Wronskian and applying integral formulas.}
    \textbf{Questions: 19, 20, 21}
    
    \item \textbf{Concept 8: Reduction of Order} 
    \textit{Using a known solution $y_1(t)$ and the standard form $p(t)$ to find a second linearly independent solution $y_2(t)$.}
    \textbf{Questions: 22, 23, 24}
    
    \item \textbf{Concept 9: Abel's Theorem and the Wronskian} 
    \textit{Calculating the Wronskian directly from the differential equation using $W(t) = C \exp(-\int p(t)dt)$.}
    \textbf{Questions: 25, 26}
    
    \item \textbf{Concept 10: Substitution Method ($t = e^x$)} 
    \textit{Transforming a Cauchy-Euler equation into a constant-coefficient equation using the substitution $x = \ln t$.}
    \textbf{Questions: 27, 28}
    
    \item \textbf{Concept 11: Error Analysis / "Find the Flaw"} 
    \textit{Identifying common pitfalls such as the $(b-a)$ trap and the standard form mandate.}
    \textbf{Questions: 29, 30}
\end{itemize}

\newpage

% ---------------------------------------------------------
\section{Problem Sets}

\subsection*{Concept 1: Existence and Uniqueness}
\textbf{1.} Determine the largest intervals for which Theorem 5 ensures the existence and uniqueness of a solution to $t(t - 2)y'' + 3ty' - 4y = t^3$. \\
\textbf{2.} Determine the largest intervals for which existence and uniqueness are guaranteed for the equation $(t^2 - 9)y'' + ty' + y = \ln(t)$. \\
\textbf{3.} Determine the largest interval for which existence and uniqueness are guaranteed for the equation $t^2 y'' + \sqrt{t - 1} y' - y = e^t$.

\subsection*{Concept 2: Homogeneous Cauchy-Euler (Distinct Real Roots)}
\textbf{4.} Find the general solution to $t^2 y'' - 2ty' + 2y = 0$ for $t > 0$. \\
\textbf{5.} Find the general solution to $2t^2 y'' + 5ty' - 2y = 0$ for $t > 0$. \\
\textbf{6.} Find the general solution to $t^2 y'' + 4ty' - 10y = 0$ for $t > 0$.

\subsection*{Concept 3: Homogeneous Cauchy-Euler (Repeated Real Roots)}
\textbf{7.} Find the general solution to $t^2 y'' - 5ty' + 9y = 0$ for $t > 0$. \\
\textbf{8.} Find the general solution to $4t^2 y'' + 8ty' + y = 0$ for $t > 0$. \\
\textbf{9.} Find the general solution to $t^2 y'' + 7ty' + 9y = 0$ for $t > 0$.

\subsection*{Concept 4: Homogeneous Cauchy-Euler (Complex Roots)}
\textbf{10.} Find the general solution to $t^2 y'' - ty' + 5y = 0$ for $t > 0$. \\
\textbf{11.} Find the general solution to $t^2 y'' + 3ty' + 10y = 0$ for $t > 0$. \\
\textbf{12.} Find the general solution to $t^2 y'' + 5ty' + 13y = 0$ for $t > 0$.

\subsection*{Concept 5: Domain Adaptation ($t < 0$)}
\textbf{13.} Find the general solution to $t^2 y'' - 3ty' + 4y = 0$ for $t < 0$. \\
\textbf{14.} Find the general solution to $t^2 y'' + 3ty' + 5y = 0$ for $t < 0$. \\
\textbf{15.} Find the general solution to $t^2 y'' + ty' - 4y = 0$ for $t < 0$.

\subsection*{Concept 6: Initial Value Problems}
\textbf{16.} Solve the initial value problem $t^2 y'' - 4ty' + 6y = 0$, with $y(1) = 1$ and $y'(1) = 0$. \\
\textbf{17.} Solve the initial value problem $t^2 y'' + 5ty' + 4y = 0$, with $y(1) = 0$ and $y'(1) = 2$. \\
\textbf{18.} Solve the initial value problem $t^2 y'' - ty' + 2y = 0$, with $y(1) = 2$ and $y'(1) = 1$.

\subsection*{Concept 7: Nonhomogeneous Equations (Variation of Parameters)}
\textbf{19.} Find the general solution to $t^2 y'' - 2ty' + 2y = t^4$ for $t > 0$. \\
\textbf{20.} Find the general solution to $t^2 y'' + ty' - y = t \ln t$ for $t > 0$. \\
\textbf{21.} Find the general solution to $t^2 y'' - 3ty' + 3y = 2t^4 \sin(t)$ for $t > 0$.

\subsection*{Concept 8: Reduction of Order}
\textbf{22.} Given the differential equation $t^2 y'' - 4ty' + 6y = 0$ for $t > 0$ and a known solution $y_1(t) = t^2$, find a second linearly independent solution using reduction of order. \\
\textbf{23.} Given the differential equation $t^2 y'' + 6ty' + 6y = 0$ for $t > 0$ and a known solution $y_1(t) = t^{-2}$, find a second linearly independent solution using reduction of order. \\
\textbf{24.} Given the differential equation $t^2 y'' - ty' + y = 0$ for $t > 0$ and a known solution $y_1(t) = t$, find a second linearly independent solution using reduction of order.

\subsection*{Concept 9: Abel's Theorem and the Wronskian}
\textbf{25.} Use Abel's formula to find the Wronskian (up to an arbitrary constant $C$) of any two solutions to $t^2 y'' + 5ty' - 3y = 0$ for $t > 0$. \\
\textbf{26.} The Wronskian $W(y_1, y_2)(t)$ of two solutions to $t^2 y'' - 2ty' + 7y = 0$ is known to be $4$ at $t = 2$. Find the value of the Wronskian at $t = 4$.

\subsection*{Concept 10: Substitution Method ($t = e^x$)}
\textbf{27.} Solve the nonhomogeneous equation $t^2 y'' - 3ty' + 4y = t^2$ for $t > 0$ by using the substitution $t = e^x$ to transform it into a constant-coefficient equation. \\
\textbf{28.} Solve the nonhomogeneous equation $t^2 y'' + ty' + y = \ln t$ for $t > 0$ by using the substitution $t = e^x$.

\subsection*{Concept 11: Error Analysis / "Find the Flaw"}
\textbf{29.} \textbf{Task:} Solve $t^2 y'' - 5ty' + 8y = 0$ for $t > 0$. \\
\textbf{Flawed Solution:} Assume $y = t^r$. The characteristic equation is $r^2 - 5r + 8 = 0$. Using the quadratic formula, $r = \frac{5 \pm \sqrt{25 - 32}}{2} = \frac{5}{2} \pm i\frac{\sqrt{7}}{2}$. The general solution is $y(t) = t^{5/2}\left(c_1 \cos\left(\frac{\sqrt{7}}{2}\ln t\right) + c_2 \sin\left(\frac{\sqrt{7}}{2}\ln t\right)\right)$. \\
\textit{Identify the error and provide the correct solution.}

\textbf{30.} \textbf{Task:} Find a particular solution to $t^2 y'' + 2ty' - 2y = t^3$ for $t > 0$. \\
\textbf{Flawed Solution:} The homogeneous roots come from $r^2 + r - 2 = 0 \implies (r+2)(r-1)=0$. So $y_1 = t$ and $y_2 = t^{-2}$. The Wronskian is $W = t(-2t^{-3}) - (1)(t^{-2}) = -3t^{-2}$. Using Variation of Parameters, $v_1 = \int \frac{-t^3(t^{-2})}{-3t^{-2}} dt = \int \frac{1}{3} t^3 dt = \frac{t^4}{12}$. $v_2 = \int \frac{t^3(t)}{-3t^{-2}} dt = \int -\frac{1}{3} t^6 dt = -\frac{t^7}{21}$. Therefore, $y_p = v_1 y_1 + v_2 y_2 = \left(\frac{t^4}{12}\right)t + \left(-\frac{t^7}{21}\right)t^{-2} = \frac{t^5}{12} - \frac{t^5}{21} = \frac{t^5}{28}$. \\
\textit{Identify the error and provide the correct solution.}

\newpage

% ---------------------------------------------------------
\section{Detailed Solutions}

\subsection*{Concept 1: Existence and Uniqueness}

\textbf{Solution 1:}
First, put the equation into standard form $y'' + p(t)y' + q(t)y = g(t)$ by dividing by $t(t - 2)$:
\[ y'' + \frac{3}{t - 2}y' - \frac{4}{t(t - 2)}y = \frac{t^2}{t - 2} \]
Identify the coefficient functions:
\[ p(t) = \frac{3}{t - 2}, \quad q(t) = \frac{-4}{t(t - 2)}, \quad g(t) = \frac{t^2}{t - 2} \]
The discontinuities occur where the denominators are zero: $t = 0$ and $t = 2$. These points divide the real number line into three open intervals of continuity:
\textbf{Final Answer:} $(-\infty, 0)$, $(0, 2)$, and $(2, \infty)$.

\textbf{Solution 2:}
Divide by $(t^2 - 9)$ to achieve standard form:
\[ y'' + \frac{t}{t^2 - 9}y' + \frac{1}{t^2 - 9}y = \frac{\ln(t)}{t^2 - 9} \]
Identify the coefficient functions:
\[ p(t) = \frac{t}{t^2 - 9}, \quad q(t) = \frac{1}{t^2 - 9}, \quad g(t) = \frac{\ln(t)}{t^2 - 9} \]
The functions $p(t)$ and $q(t)$ are discontinuous at $t = \pm 3$. The function $g(t)$ requires $t > 0$ due to the natural logarithm and is discontinuous at $t = 3$. The valid domain is $t > 0$, split by the discontinuity at $t = 3$.
\textbf{Final Answer:} $(0, 3)$ and $(3, \infty)$.

\textbf{Solution 3:}
Divide by $t^2$ to achieve standard form:
\[ y'' + \frac{\sqrt{t - 1}}{t^2}y' - \frac{1}{t^2}y = \frac{e^t}{t^2} \]
Identify the coefficient functions:
\[ p(t) = \frac{\sqrt{t - 1}}{t^2}, \quad q(t) = \frac{-1}{t^2}, \quad g(t) = \frac{e^t}{t^2} \]
The function $p(t)$ requires $t - 1 \ge 0 \implies t \ge 1$. The denominator $t^2$ introduces a discontinuity at $t = 0$, but this is already outside the domain of $p(t)$. The intersection of their continuous domains is $(1, \infty)$.
\textbf{Final Answer:} $(1, \infty)$.

\subsection*{Concept 2: Homogeneous Cauchy-Euler (Distinct Real Roots)}

\textbf{Solution 4:}
This is a Cauchy-Euler equation with $a=1, b=-2, c=2$. The characteristic equation is $ar^2 + (b-a)r + c = 0$:
\[ (1)r^2 + (-2 - 1)r + 2 = 0 \implies r^2 - 3r + 2 = 0 \]
Factoring yields $(r - 1)(r - 2) = 0$. The roots are $r_1 = 1$ and $r_2 = 2$.
\textbf{Final Answer:} $y(t) = c_1 t + c_2 t^2$.

\textbf{Solution 5:}
Here, $a=2, b=5, c=-2$. The characteristic equation is:
\[ 2r^2 + (5 - 2)r - 2 = 0 \implies 2r^2 + 3r - 2 = 0 \]
Factoring yields $(2r - 1)(r + 2) = 0$. The roots are $r_1 = 1/2$ and $r_2 = -2$.
\textbf{Final Answer:} $y(t) = c_1 t^{1/2} + c_2 t^{-2}$.

\textbf{Solution 6:}
Here, $a=1, b=4, c=-10$. The characteristic equation is:
\[ r^2 + (4 - 1)r - 10 = 0 \implies r^2 + 3r - 10 = 0 \]
Factoring yields $(r + 5)(r - 2) = 0$. The roots are $r_1 = -5$ and $r_2 = 2$.
\textbf{Final Answer:} $y(t) = c_1 t^{-5} + c_2 t^2$.

\subsection*{Concept 3: Homogeneous Cauchy-Euler (Repeated Real Roots)}

\textbf{Solution 7:}
Here, $a=1, b=-5, c=9$. The characteristic equation is:
\[ r^2 + (-5 - 1)r + 9 = 0 \implies r^2 - 6r + 9 = 0 \]
Factoring yields $(r - 3)^2 = 0$. We have a repeated root $r = 3$. The second linearly independent solution requires multiplying by $\ln t$.
\textbf{Final Answer:} $y(t) = c_1 t^3 + c_2 t^3 \ln t$.

\textbf{Solution 8:}
Here, $a=4, b=8, c=1$. The characteristic equation is:
\[ 4r^2 + (8 - 4)r + 1 = 0 \implies 4r^2 + 4r + 1 = 0 \]
Factoring yields $(2r + 1)^2 = 0$. We have a repeated root $r = -1/2$.
\textbf{Final Answer:} $y(t) = c_1 t^{-1/2} + c_2 t^{-1/2} \ln t$.

\textbf{Solution 9:}
Here, $a=1, b=7, c=9$. The characteristic equation is:
\[ r^2 + (7 - 1)r + 9 = 0 \implies r^2 + 6r + 9 = 0 \]
Factoring yields $(r + 3)^2 = 0$. We have a repeated root $r = -3$.
\textbf{Final Answer:} $y(t) = c_1 t^{-3} + c_2 t^{-3} \ln t$.

\subsection*{Concept 4: Homogeneous Cauchy-Euler (Complex Roots)}

\textbf{Solution 10:}
Here, $a=1, b=-1, c=5$. The characteristic equation is:
\[ r^2 + (-1 - 1)r + 5 = 0 \implies r^2 - 2r + 5 = 0 \]
Using the quadratic formula: $r = \frac{2 \pm \sqrt{4 - 20}}{2} = 1 \pm 2i$. Here $\alpha = 1, \beta = 2$.
\textbf{Final Answer:} $y(t) = t(c_1 \cos(2\ln t) + c_2 \sin(2\ln t))$.

\textbf{Solution 11:}
Here, $a=1, b=3, c=10$. The characteristic equation is:
\[ r^2 + (3 - 1)r + 10 = 0 \implies r^2 + 2r + 10 = 0 \]
Using the quadratic formula: $r = \frac{-2 \pm \sqrt{4 - 40}}{2} = -1 \pm 3i$. Here $\alpha = -1, \beta = 3$.
\textbf{Final Answer:} $y(t) = t^{-1}(c_1 \cos(3\ln t) + c_2 \sin(3\ln t))$.

\textbf{Solution 12:}
Here, $a=1, b=5, c=13$. The characteristic equation is:
\[ r^2 + (5 - 1)r + 13 = 0 \implies r^2 + 4r + 13 = 0 \]
Using the quadratic formula: $r = \frac{-4 \pm \sqrt{16 - 52}}{2} = -2 \pm 3i$. Here $\alpha = -2, \beta = 3$.
\textbf{Final Answer:} $y(t) = t^{-2}(c_1 \cos(3\ln t) + c_2 \sin(3\ln t))$.

\subsection*{Concept 5: Domain Adaptation ($t < 0$)}

\textbf{Solution 13:}
Here, $a=1, b=-3, c=4$. The characteristic equation is:
\[ r^2 + (-3 - 1)r + 4 = 0 \implies r^2 - 4r + 4 = 0 \implies (r - 2)^2 = 0 \]
The root is $r = 2$ (repeated). Because $t < 0$, $\ln t$ is undefined. We must use $\ln(-t)$.
\textbf{Final Answer:} $y(t) = c_1 t^2 + c_2 t^2 \ln(-t)$.

\textbf{Solution 14:}
Here, $a=1, b=3, c=5$. The characteristic equation is:
\[ r^2 + (3 - 1)r + 5 = 0 \implies r^2 + 2r + 5 = 0 \]
Roots are $r = -1 \pm 2i$. Because $t < 0$, we use $-t$ inside the logarithm and for the base power (though $(-t)^{-1}$ and $t^{-1}$ only differ by a negative sign which is absorbed by constants, standard convention is to use $|t|$).
\textbf{Final Answer:} $y(t) = (-t)^{-1}(c_1 \cos(2\ln(-t)) + c_2 \sin(2\ln(-t)))$.

\textbf{Solution 15:}
Here, $a=1, b=1, c=-4$. The characteristic equation is:
\[ r^2 + (1 - 1)r - 4 = 0 \implies r^2 - 4 = 0 \implies r = \pm 2 \]
The roots are real and distinct. No logarithms are present.
\textbf{Final Answer:} $y(t) = c_1 t^2 + c_2 t^{-2}$.

\subsection*{Concept 6: Initial Value Problems}

\textbf{Solution 16:}
Characteristic equation: $r^2 + (-4 - 1)r + 6 = 0 \implies r^2 - 5r + 6 = 0 \implies (r-2)(r-3)=0$.
General solution: $y(t) = c_1 t^2 + c_2 t^3$.
Derivative: $y'(t) = 2c_1 t + 3c_2 t^2$.
Apply initial conditions at $t=1$:
$y(1) = c_1 + c_2 = 1$
$y'(1) = 2c_1 + 3c_2 = 0$
From the first equation, $c_1 = 1 - c_2$. Substitute into the second: $2(1 - c_2) + 3c_2 = 0 \implies 2 + c_2 = 0 \implies c_2 = -2$. Then $c_1 = 3$.
\textbf{Final Answer:} $y(t) = 3t^2 - 2t^3$.

\textbf{Solution 17:}
Characteristic equation: $r^2 + (5 - 1)r + 4 = 0 \implies r^2 + 4r + 4 = 0 \implies (r+2)^2 = 0$.
General solution: $y(t) = c_1 t^{-2} + c_2 t^{-2} \ln t$.
Derivative: $y'(t) = -2c_1 t^{-3} + c_2(t^{-3} - 2t^{-3} \ln t)$.
Apply initial conditions at $t=1$:
$y(1) = c_1(1) + c_2(0) = 0 \implies c_1 = 0$.
$y'(1) = -2(0)(1) + c_2(1 - 0) = 2 \implies c_2 = 2$.
\textbf{Final Answer:} $y(t) = 2t^{-2} \ln t$.

\textbf{Solution 18:}
Characteristic equation: $r^2 + (-1 - 1)r + 2 = 0 \implies r^2 - 2r + 2 = 0 \implies r = 1 \pm i$.
General solution: $y(t) = t(c_1 \cos(\ln t) + c_2 \sin(\ln t))$.
Derivative: $y'(t) = (c_1 \cos(\ln t) + c_2 \sin(\ln t)) + t(-c_1 \sin(\ln t)\frac{1}{t} + c_2 \cos(\ln t)\frac{1}{t}) = (c_1 + c_2)\cos(\ln t) + (c_2 - c_1)\sin(\ln t)$.
Apply initial conditions at $t=1$:
$y(1) = 1(c_1(1) + c_2(0)) = 2 \implies c_1 = 2$.
$y'(1) = (c_1 + c_2)(1) + (c_2 - c_1)(0) = c_1 + c_2 = 1 \implies 2 + c_2 = 1 \implies c_2 = -1$.
\textbf{Final Answer:} $y(t) = t(2\cos(\ln t) - \sin(\ln t))$.

\subsection*{Concept 7: Nonhomogeneous Equations (Variation of Parameters)}

\textbf{Solution 19:}
Step 1: Homogeneous solution. $r^2 - 3r + 2 = 0 \implies r=1, 2$. $y_1 = t, y_2 = t^2$.
Step 2: Standard form. Divide by $t^2$: $y'' - \frac{2}{t}y' + \frac{2}{t^2}y = t^2$. So $g(t) = t^2$.
Step 3: Wronskian. $W = y_1 y_2' - y_1' y_2 = t(2t) - (1)t^2 = t^2$.
Step 4: Variation of Parameters.
$v_1 = \int \frac{-g(t)y_2(t)}{W} dt = \int \frac{-t^2(t^2)}{t^2} dt = \int -t^2 dt = -\frac{t^3}{3}$.
$v_2 = \int \frac{g(t)y_1(t)}{W} dt = \int \frac{t^2(t)}{t^2} dt = \int t dt = \frac{t^2}{2}$.
$y_p = v_1 y_1 + v_2 y_2 = \left(-\frac{t^3}{3}\right)t + \left(\frac{t^2}{2}\right)t^2 = -\frac{t^4}{3} + \frac{t^4}{2} = \frac{t^4}{6}$.
\textbf{Final Answer:} $y(t) = c_1 t + c_2 t^2 + \frac{t^4}{6}$.

\textbf{Solution 20:}
Step 1: Homogeneous solution. $r^2 + (1-1)r - 1 = 0 \implies r^2 - 1 = 0 \implies r = \pm 1$. $y_1 = t, y_2 = t^{-1}$.
Step 2: Standard form. Divide by $t^2$: $g(t) = \frac{t \ln t}{t^2} = t^{-1} \ln t$.
Step 3: Wronskian. $W = t(-t^{-2}) - (1)t^{-1} = -2t^{-1}$.
Step 4: Variation of Parameters.
$v_1 = \int \frac{-(t^{-1} \ln t)(t^{-1})}{-2t^{-1}} dt = \int \frac{1}{2} t^{-1} \ln t dt = \frac{1}{4} (\ln t)^2$.
$v_2 = \int \frac{(t^{-1} \ln t)(t)}{-2t^{-1}} dt = \int -\frac{1}{2} \ln t dt = -\frac{1}{2}(t \ln t - t)$.
$y_p = v_1 y_1 + v_2 y_2 = \left(\frac{1}{4} (\ln t)^2\right)t + \left(-\frac{1}{2} t \ln t + \frac{1}{2} t\right)t^{-1} = \frac{1}{4} t (\ln t)^2 - \frac{1}{2} \ln t + \frac{1}{2}$.
\textbf{Final Answer:} $y(t) = c_1 t + c_2 t^{-1} + \frac{1}{4} t (\ln t)^2 - \frac{1}{2} \ln t$. (Note: the constant $1/2$ can be absorbed into the homogeneous solution if desired, but keeping it is perfectly valid).

\textbf{Solution 21:}
Step 1: Homogeneous solution. $r^2 - 4r + 3 = 0 \implies r=1, 3$. $y_1 = t, y_2 = t^3$.
Step 2: Standard form. Divide by $t^2$: $g(t) = 2t^2 \sin(t)$.
Step 3: Wronskian. $W = t(3t^2) - (1)t^3 = 2t^3$.
Step 4: Variation of Parameters.
$v_1 = \int \frac{-2t^2 \sin(t) (t^3)}{2t^3} dt = \int -t^2 \sin(t) dt$. Using integration by parts twice yields $v_1 = t^2 \cos(t) - 2t \sin(t) - 2\cos(t)$.
$v_2 = \int \frac{2t^2 \sin(t) (t)}{2t^3} dt = \int \sin(t) dt = -\cos(t)$.
$y_p = (t^2 \cos(t) - 2t \sin(t) - 2\cos(t))t + (-\cos(t))t^3 = -2t^2 \sin(t) - 2t \cos(t)$.
\textbf{Final Answer:} $y(t) = c_1 t + c_2 t^3 - 2t^2 \sin(t) - 2t \cos(t)$.

\subsection*{Concept 8: Reduction of Order}

\textbf{Solution 22:}
Put the equation in standard form: $y'' - \frac{4}{t}y' + \frac{6}{t^2}y = 0$. Identify $p(t) = -\frac{4}{t}$.
Calculate the integrating factor: $e^{-\int p(t)dt} = e^{\int \frac{4}{t} dt} = e^{4\ln t} = t^4$.
Apply the reduction of order formula:
$y_2(t) = y_1(t) \int \frac{e^{-\int p(t)dt}}{y_1(t)^2} dt = t^2 \int \frac{t^4}{(t^2)^2} dt = t^2 \int 1 dt = t^3$.
\textbf{Final Answer:} $y_2(t) = t^3$.

\textbf{Solution 23:}
Put the equation in standard form: $y'' + \frac{6}{t}y' + \frac{6}{t^2}y = 0$. Identify $p(t) = \frac{6}{t}$.
Calculate the integrating factor: $e^{-\int p(t)dt} = e^{-\int \frac{6}{t} dt} = e^{-6\ln t} = t^{-6}$.
Apply the reduction of order formula:
$y_2(t) = t^{-2} \int \frac{t^{-6}}{(t^{-2})^2} dt = t^{-2} \int \frac{t^{-6}}{t^{-4}} dt = t^{-2} \int t^{-2} dt = t^{-2}(-t^{-1}) = -t^{-3}$.
Dropping the constant multiple $-1$, we get $t^{-3}$.
\textbf{Final Answer:} $y_2(t) = t^{-3}$.

\textbf{Solution 24:}
Put the equation in standard form: $y'' - \frac{1}{t}y' + \frac{1}{t^2}y = 0$. Identify $p(t) = -\frac{1}{t}$.
Calculate the integrating factor: $e^{-\int p(t)dt} = e^{\int \frac{1}{t} dt} = e^{\ln t} = t$.
Apply the reduction of order formula:
$y_2(t) = t \int \frac{t}{(t)^2} dt = t \int \frac{1}{t} dt = t \ln t$.
\textbf{Final Answer:} $y_2(t) = t \ln t$.

\subsection*{Concept 9: Abel's Theorem and the Wronskian}

\textbf{Solution 25:}
First, put the equation in standard form by dividing by $t^2$: $y'' + \frac{5}{t}y' - \frac{3}{t^2}y = 0$.
Identify $p(t) = \frac{5}{t}$.
According to Abel's formula, $W(t) = C \exp\left(-\int p(t)dt\right)$.
$W(t) = C \exp\left(-\int \frac{5}{t} dt\right) = C \exp(-5\ln t) = C t^{-5}$.
\textbf{Final Answer:} $W(t) = C t^{-5}$.

\textbf{Solution 26:}
Put the equation in standard form: $y'' - \frac{2}{t}y' + \frac{7}{t^2}y = 0$. Identify $p(t) = -\frac{2}{t}$.
Using Abel's formula: $W(t) = C \exp\left(-\int -\frac{2}{t} dt\right) = C \exp(2\ln t) = C t^2$.
We are given $W(2) = 4$. Substitute $t = 2$:
$C(2)^2 = 4 \implies 4C = 4 \implies C = 1$.
Thus, the specific Wronskian is $W(t) = t^2$.
Evaluate at $t = 4$: $W(4) = (4)^2 = 16$.
\textbf{Final Answer:} $W(4) = 16$.

\subsection*{Concept 10: Substitution Method ($t = e^x$)}

\textbf{Solution 27:}
Let $t = e^x$, so $x = \ln t$. The derivatives transform as $t \frac{dy}{dt} = \frac{dy}{dx}$ and $t^2 \frac{d^2y}{dt^2} = \frac{d^2y}{dx^2} - \frac{dy}{dx}$.
Substitute into the original equation:
$\left(\frac{d^2y}{dx^2} - \frac{dy}{dx}\right) - 3\left(\frac{dy}{dx}\right) + 4y = (e^x)^2$
$\frac{d^2y}{dx^2} - 4\frac{dy}{dx} + 4y = e^{2x}$
This is a constant-coefficient equation. The homogeneous roots are $r^2 - 4r + 4 = 0 \implies (r-2)^2 = 0$. So $y_h(x) = c_1 e^{2x} + c_2 x e^{2x}$.
Using Undetermined Coefficients, guess $y_p(x) = A x^2 e^{2x}$ (since $e^{2x}$ and $x e^{2x}$ are in the homogeneous solution).
$y_p' = 2Ax e^{2x} + 2Ax^2 e^{2x}$
$y_p'' = 2A e^{2x} + 8Ax e^{2x} + 4Ax^2 e^{2x}$
Substitute into the ODE:
$(2A + 8Ax + 4Ax^2 - 8Ax - 8Ax^2 + 4Ax^2)e^{2x} = e^{2x} \implies 2A = 1 \implies A = 1/2$.
So $y(x) = c_1 e^{2x} + c_2 x e^{2x} + \frac{1}{2} x^2 e^{2x}$.
Substitute back $x = \ln t$ and $e^{2x} = t^2$:
\textbf{Final Answer:} $y(t) = c_1 t^2 + c_2 t^2 \ln t + \frac{1}{2} t^2 (\ln t)^2$.

\textbf{Solution 28:}
Let $t = e^x$, so $x = \ln t$. Substitute into the original equation:
$\left(\frac{d^2y}{dx^2} - \frac{dy}{dx}\right) + \left(\frac{dy}{dx}\right) + y = x$
$\frac{d^2y}{dx^2} + y = x$
The homogeneous roots are $r^2 + 1 = 0 \implies r = \pm i$. So $y_h(x) = c_1 \cos x + c_2 \sin x$.
Using Undetermined Coefficients, guess $y_p(x) = Ax + B$.
$y_p'' = 0$. Substitute into the ODE: $0 + Ax + B = x \implies A = 1, B = 0$.
So $y(x) = c_1 \cos x + c_2 \sin x + x$.
Substitute back $x = \ln t$:
\textbf{Final Answer:} $y(t) = c_1 \cos(\ln t) + c_2 \sin(\ln t) + \ln t$.

\subsection*{Concept 11: Error Analysis / "Find the Flaw"}

\textbf{Solution 29:}
\textbf{Error Analysis:} The flawed solution falls into the $(b-a)$ trap. The characteristic equation for a Cauchy-Euler equation $at^2y'' + bty' + cy = 0$ is $ar^2 + (b - a)r + c = 0$, not $ar^2 + br + c = 0$. The student forgot to subtract $a$ from $b$ for the middle term.
\textbf{Correct Step:} The correct characteristic equation is $r^2 + (-5 - 1)r + 8 = 0 \implies r^2 - 6r + 8 = 0$.
Factoring yields $(r - 2)(r - 4) = 0$. The roots are real and distinct: $r = 2, 4$.
\textbf{Correct Final Solution:} $y(t) = c_1 t^2 + c_2 t^4$.

\textbf{Solution 30:}
\textbf{Error Analysis:} The flawed solution violates the Standard Form Mandate. The function $g(t)$ in the Variation of Parameters formula must be taken from the standard form $y'' + p(t)y' + q(t)y = g(t)$. The student used $g(t) = t^3$ from the original equation instead of dividing by the leading coefficient $t^2$.
\textbf{Correct Step:} Divide the entire equation by $t^2$ to get $y'' + \frac{2}{t}y' - \frac{2}{t^2}y = t$. Thus, the correct $g(t) = t$.
Recalculate the integrals:
$v_1 = \int \frac{-t(t^{-2})}{-3t^{-2}} dt = \int \frac{1}{3} t dt = \frac{t^2}{6}$.
$v_2 = \int \frac{t(t)}{-3t^{-2}} dt = \int -\frac{1}{3} t^4 dt = -\frac{t^5}{15}$.
$y_p = v_1 y_1 + v_2 y_2 = \left(\frac{t^2}{6}\right)t + \left(-\frac{t^5}{15}\right)t^{-2} = \frac{t^3}{6} - \frac{t^3}{15} = \frac{5t^3 - 2t^3}{30} = \frac{3t^3}{30} = \frac{t^3}{10}$.
\textbf{Correct Final Solution:} $y_p(t) = \frac{t^3}{10}$.

\newpage

% ---------------------------------------------------------
\section{Concept Checklist \& Question Mapping (Appendix)}
\textit{As requested, here is the reference mapping to easily locate problems by concept.}

\begin{itemize}[label=\checkmark]
    \item \textbf{Existence and Uniqueness Analysis:} Questions 1, 2, 3
    \item \textbf{Homogeneous Cauchy-Euler (Distinct Real Roots):} Questions 4, 5, 6
    \item \textbf{Homogeneous Cauchy-Euler (Repeated Real Roots):} Questions 7, 8, 9
    \item \textbf{Homogeneous Cauchy-Euler (Complex Roots):} Questions 10, 11, 12
    \item \textbf{Domain Adaptation ($t < 0$):} Questions 13, 14, 15
    \item \textbf{Initial Value Problems (IVPs):} Questions 16, 17, 18
    \item \textbf{Nonhomogeneous Equations (Variation of Parameters):} Questions 19, 20, 21
    \item \textbf{Reduction of Order:} Questions 22, 23, 24
    \item \textbf{Abel's Theorem and the Wronskian:} Questions 25, 26
    \item \textbf{Substitution Method ($t = e^x$):} Questions 27, 28
    \item \textbf{Error Analysis / "Find the Flaw":} Questions 29, 30
\end{itemize}

\end{document}